\section{XOR：隐藏表示怎样改变可分性}
\label{sec:nn-xor}
异或（XOR）在两个输入不同的时候输出1，相同时输出0。
四个输入$(0,0),(1,0),(0,1),(1,1)$的目标为$0,1,1,0$；
它们不能由原平面中的一条直线严格分开。

\begin{figure}[H]
\centering\begin{tikzpicture}
\begin{axis}[width=68mm,height=58mm,xmin=-.3,xmax=1.3,ymin=-.3,ymax=1.3,
 xlabel={$x_1$},ylabel={$x_2$},xtick={0,1},ytick={0,1},title={原输入：一条直线无法分开}]
\addplot[only marks,mark=*,Teal,mark size=3pt] coordinates {(1,0)(0,1)};
\addplot[only marks,mark=square*,Amber,mark size=3pt] coordinates {(0,0)(1,1)};
\end{axis}
\begin{axis}[at={(79mm,0)},anchor=south west,width=68mm,height=58mm,xmin=-.3,xmax=1.3,ymin=-.3,ymax=1.3,
 xlabel={$h_1=\ReLU(x_1-x_2)$},ylabel={$h_2=\ReLU(x_2-x_1)$},xtick={0,1},ytick={0,1},title={新表示：按$h_1+h_2$分开}]
\addplot[Quiet,dashed,domain=-.2:.7]{.5-x};
\addplot[only marks,mark=*,Teal,mark size=3pt] coordinates {(1,0)(0,1)};
\addplot[only marks,mark=square*,Amber,mark size=3pt] coordinates {(0,0)};
\end{axis}
\end{tikzpicture}

\caption{青色圆点为正类，橙色方点为负类；两负类在隐藏表示中重合到原点。}
\end{figure}
\readfigure{设$h_1=\ReLU(x_1-x_2)$、$h_2=\ReLU(x_2-x_1)$，
则正类映为$(1,0)$或$(0,1)$，负类均为$(0,0)$。
再按$h_1+h_2>0.5$即可分类。这里显式构造权重，说明表示能力，不等于保证训练自动找到它。}

若所有层都只有仿射变换，$W_2(W_1x+b_1)+b_2$
仍是一个仿射变换；单纯叠层没有引入这里所需的非线性表达。
中间的非线性使组合获得不同形状的决策边界。

深度可让某些具有组合结构的函数更高效地表示，并不意味着越深越好。
常见通用逼近结论有函数、区域与网络宽度等条件，也不保证优化、样本需求或测试泛化。
XOR只需很小的隐藏表示，不需要很深的网络。

\quick{先问结构能否表达，再问优化能否学到，最后问新数据是否有效；三个问题不能相互替代。}

